\chapter{Giới thiệu và Mục tiêu}
\label{ch:gioithieu}

\section{Bối cảnh}
Phần mềm thanh toán là nơi một lỗi logic nhỏ chuyển thành thiệt hại tài chính trực tiếp:
giao hàng khi tiền chưa thực sự được thu, ghi nhận thu tiền hai lần khi cổng thanh toán gửi
lại webhook, hoặc chấp nhận một webhook giả mạo nhờ một ``mã bí mật'' bị bỏ quên trong mã
nguồn. Các lỗi này thuộc nhóm \cwe{840}{Business Logic Errors}, \cwe{798}{Hard-coded
Credentials} và \cwe{667}{Improper Locking} đã học ở Chương~1. Bên cạnh đó, các hàm xử lý
dữ liệu từ cổng thanh toán (phân tích chữ ký webhook, cộng dồn số tiền) còn chịu các lớp lỗi
bộ nhớ và số học của Chương~2--3.

Báo cáo áp dụng tuần tự kỹ thuật của cả bốn chương lên \emph{một} hệ thống — cổng thanh
toán đơn hàng thương mại điện tử \code{payment\_gateway} — theo hướng Chủ đề~5 (quy trình
end-to-end) của Hướng dẫn Bài tập lớn. Hệ thống được đặc tả hai lần: \textbf{v0} là thiết kế
ban đầu còn lỗi, \textbf{v1} là thiết kế đã sửa theo kết quả kiểm chứng.

\section{Mục tiêu cụ thể}
\begin{enumerate}
  \item \textbf{Ch.1 — Mô hình hóa:} dựng cấu trúc Kripke $K=(S,S_0,R,L)$ cho vòng đời thanh
        toán và khóa sổ cái; kiểm chứng 8 công thức CTL (an toàn, sống còn, khả đạt) và 5 cặp
        loại trừ lẫn nhau, có phản ví dụ cho mọi vi phạm.
  \item \textbf{Ch.3 — BMC:} kiểm chứng lại các thuộc tính an toàn bằng Bounded Model Checking
        với Z3, xác định completeness threshold để kết luận ``an toàn'' là đầy đủ.
  \item \textbf{Ch.2 — Bộ nhớ:} phát hiện và vá 4 lỗi bộ nhớ C (tràn bộ đệm, rò rỉ, double
        free, use-after-free), xác nhận trước/sau bằng static scanner và AddressSanitizer.
  \item \textbf{Ch.3 — SMT:} chứng minh khả năng tràn int32 của 6 bộ đếm/tổng tiền, và chứng
        minh tiền điều kiện của bản vá loại bỏ hoàn toàn tràn số (UNSAT).
  \item \textbf{Ch.4 — Fuzzing:} so sánh định lượng black-box, coverage-guided và concolic
        fuzzing với cùng ngân sách 2\,000 lần thử trên 3 hàm Python.
\end{enumerate}

\section{Phạm vi}
\textbf{Trong phạm vi:} một đơn hàng, thanh toán toàn phần qua một cổng; hoàn tiền toàn phần;
khóa sổ cái giữa hai worker; 9 hàm C/Python sinh ra từ đặc tả (khoảng 340 dòng mã C cho hai
phiên bản); kiểm chứng trên mô hình và trên mã sinh ra.

\textbf{Ngoài phạm vi:} hoàn tiền một phần, nhiều cổng thanh toán, tích hợp ngân hàng thật,
hiệu năng (các yêu cầu phi chức năng chỉ được tài liệu hóa, không kiểm chứng tự động).

\begin{ghichu}[Tuyên bố về lỗi cài đặt có chủ đích]
Các lỗi trong đặc tả v0 (cạnh chuyển trạng thái sai trong mô hình, cột \code{Loi} của sheet
\code{HamSinhMa}) được \textbf{cài có chủ đích} để làm đối tượng cho công cụ phát hiện, đúng như
Hướng dẫn Bài tập lớn cho phép. Mọi kết luận ``phát hiện được'' trong báo cáo đều dựa trên
\emph{kết quả chạy công cụ}: công cụ gộp CWE chỉ ghi nhận một lỗ hổng khi có bằng chứng
(finding của scanner, log sanitizer, model của Z3, input gây lỗi của fuzzer, phản ví dụ
CTL/BMC); lỗi đã khai báo mà không công cụ nào tái hiện được sẽ bị đánh dấu
\emph{declared-only}. Trong lần chạy cuối, không có lỗi nào thuộc loại này.
\end{ghichu}

\section{Cấu trúc báo cáo}
Chương~\ref{ch:hethong} mô tả hệ thống mục tiêu; Chương~\ref{ch:phuongphap} trình bày
phương pháp theo từng chương học phần; Chương~\ref{ch:ketqua} trình bày kết quả kèm bằng
chứng; Chương~\ref{ch:thaoluan} thảo luận hạn chế; Chương~\ref{ch:ketluan} kết luận.
